\chapter{Summary, Conclusion, and Recommendations}
\label{ch:conclusion}

\section{Introduction}

This final chapter draws the project documentation to a close. It summarises
the study reported in Chapters~\ref{ch:introduction}--\ref{ch:implementation},
assesses the extent to which the specific objectives set out in
Section~\ref{sec:aim-objectives} were achieved by the delivered system,
states the overall conclusion that follows from that assessment, and offers
recommendations -- both to whoever deploys and operates EcoTrace going
forward, and to any team that continues its development. In keeping with the
approach taken throughout Chapter~\ref{ch:implementation} -- most notably
the module-status audit (Table~\ref{tab:module-status}) and the testing
results of Section~\ref{sec:testing} -- this chapter reports achievement
honestly against what was actually built and tested, rather than against
what was originally proposed. Non-Functional Requirement NFR-ETH-9
(Table~\ref{tab:nfr-eth}, Chapter~\ref{ch:analysis-design}) explicitly
requires the system and its documentation to distinguish implemented
functionality from planned functionality; this chapter applies that same
requirement to the project as a whole.

\section{Summary of the Study}

Electronic waste is one of the fastest-growing waste streams in the world,
and Sierra Leone, like many developing countries, lacks a coordinated
digital system for registering, collecting, tracking, and recovering value
from discarded electrical and electronic equipment (Chapter~\ref{ch:introduction}).
This study set out to design, develop, and implement EcoTrace, a
cross-platform electronic waste recovery and circular economy management
system intended to connect households and businesses that generate e-waste
with the collectors, technicians, recyclers, and administrators who process
it, while giving environmental officers and administrators the visibility
needed for monitoring and reporting.

Chapter~\ref{ch:literature} reviewed the conceptual foundations of
electronic waste management and the circular economy, surveyed the
technologies available for building a system of this kind, and examined
comparable existing systems, identifying the gaps that EcoTrace was intended
to fill. Chapter~\ref{ch:analysis-design} analysed the existing (manual and
fragmented) approach to e-waste handling, specified the proposed system's
functional and non-functional requirements as roughly 200 and 90 individual
requirement statements respectively, and presented its architecture, use
cases, data flows, sequence diagrams, database design, and user-interface
design. Chapter~\ref{ch:implementation} reported how the system was actually
built -- a Flutter client for Android, Web, and Windows; a TypeScript/Express
backend on Firebase Cloud Functions; Cloud Firestore as the primary
database; and a Monime mobile-money payment integration that was added
during development and is not part of the original proposal -- and reported
the results of the testing performed on it, including one real, live-money
end-to-end payment transaction (test case T-06, Table~\ref{tab:test-cases}).

Taken together, the four preceding chapters document a system that was
carried from a stated problem, through a formal analysis and design, to a
working, partially-tested implementation. The remainder of this chapter
evaluates that outcome against the objectives that were set for it at the
outset.

\section{Achievement of the Specific Objectives}
\label{sec:achievement-objectives}

Section~\ref{sec:aim-objectives} of Chapter~\ref{ch:introduction} set out
ten specific objectives for this study. Table~\ref{tab:objective-achievement}
assesses each objective against the implementation and testing evidence
reported in Chapter~\ref{ch:implementation}. Three status values are used,
consistent with the implementation-status classification introduced in
Chapter~\ref{ch:implementation} (Table~\ref{tab:module-status}):
\textbf{Achieved} (built, integrated, and observed working, though not
necessarily exhaustively tested); \textbf{Partially Achieved} (a working
implementation exists but departs from the objective as originally worded,
or covers only part of it); and \textbf{Not Achieved} (the objective, as
originally worded, was not delivered).

\begin{longtable}{@{}p{0.6cm}p{6.6cm}p{2.4cm}p{4.3cm}@{}}
\caption{Achievement of the Specific Objectives of the Study} \label{tab:objective-achievement} \\
\toprule \textbf{\#} & \textbf{Objective (as stated in Section~\ref{sec:aim-objectives})} & \textbf{Status} & \textbf{Evidence / Remarks} \\ \midrule
\endfirsthead
\toprule \textbf{\#} & \textbf{Objective} & \textbf{Status} & \textbf{Evidence / Remarks} \\ \midrule
\endhead
1 & Design and develop a secure cross-platform electronic waste management system for Android, Web, and Windows. & Achieved & A single Flutter codebase builds and runs for all three targets; \texttt{flutter build web} and \texttt{flutter build apk} both succeeded with the payment module included (T-15, Table~\ref{tab:test-cases}). \\
2 & Develop a centralized cloud-based platform for registering, storing, managing, and tracking e-waste information using Cloud Firestore. & Achieved & Eighteen Firestore collections are implemented and in active use (Table~\ref{tab:collections-summary}, Chapter~\ref{ch:analysis-design}; database implementation, Chapter~\ref{ch:implementation}). \\
3 & Implement secure user authentication and role-based access control for different categories of system users using Firebase Authentication. & Achieved & Firebase Authentication is live for email/password sign-in across all roles, backed by custom claims and Firestore security rules that were themselves the subject of a real bug fix during testing (T-05/T-05b, Table~\ref{tab:test-cases}); social-provider sign-in and multi-factor enrolment remain outstanding (Table~\ref{tab:module-status}). \\
4 & Integrate artificial intelligence to support automatic classification of electronic waste from uploaded images. & Partially Achieved & A pluggable classification interface was implemented and is exercised through manual and assisted review workflows, but the assisted path uses a heuristic fallback classifier rather than a trained PyTorch model as originally specified; see Section~\ref{sec:classification-module} for the reasoning behind this substitution. \\
5 & Develop an intelligent e-waste pickup scheduling and collection management module for households and businesses. & Achieved & Pickup drafting, server-side fee estimation, Cloudinary photo capture, GPS capture, scheduling, cancellation/rescheduling, and rating are implemented and tested; payment is now a required gate before a pickup can be submitted (Section~\ref{sec:pickup-module}). \\
6 & Integrate Google Maps Platform to support location services, display collection locations, visualize routes, and improve collection operations. & Not Achieved & Location capture and route generation use the \texttt{geolocator} plugin and haversine-distance calculations rather than the Google Maps Platform; no interactive map or visual route display was integrated. This is the clearest departure from the original proposal and is discussed in the conformance note at the start of Chapter~\ref{ch:implementation}. \\
7 & Develop modules for e-waste transportation, repair, refurbishment, recycling, and lifecycle tracking to promote circular economy practices. & Achieved & Collection/logistics, repair, recycling, resource-recovery, and reverse-logistics modules are all implemented and API-backed (Table~\ref{tab:module-status}). \\
8 & Implement real-time notification services using Firebase Cloud Messaging to keep users informed of requests, status updates, and system activities. & Partially Achieved & In-app and event-driven notifications are implemented via a dedicated notification microservice capable of fanning out to FCM, SMS, and email; production configuration of those external providers had not been completed at the time of writing (Table~\ref{tab:module-status}). \\
9 & Generate administrative reports and analytics to support environmental monitoring, operational planning, and informed decision-making. & Partially Achieved & Role-specific dashboards and underlying analytics queries are implemented and were exercised during testing; production PDF/XLSX export and large-dataset performance testing remain outstanding (Table~\ref{tab:module-status}). \\
10 & Evaluate the effectiveness, usability, and performance of the developed EcoTrace system in improving e-waste management processes. & Partially Achieved & The system was evaluated through structured functional and security testing (Section~\ref{sec:testing}), including one real, live-money payment transaction, but no formal usability study or field evaluation with end users was conducted; this is identified as a priority for future work in Section~\ref{sec:future-work}. \\
\bottomrule
\end{longtable}

Of the ten objectives, five were fully achieved, four were partially
achieved, and one -- the integration of the Google Maps Platform -- was not
achieved as originally specified. The single largest unplanned addition to
the system, the Monime mobile-money payment module
(Section~\ref{sec:payment-module}), was not among the original objectives
at all, but was implemented, integrated into the pickup workflow as a
required gate, and verified with a real transaction, making it the most
thoroughly tested part of the delivered system.

\section{Contributions of the Study}

Notwithstanding the partial achievement of some objectives, this study
makes the following contributions:

\begin{enumerate}[leftmargin=1.4cm]
  \item It delivers a working, cross-platform e-waste management system
    covering the full lifecycle from household/business pickup request
    through collection, repair, recycling, and resource recovery, built on
    a modern cloud-native stack (Flutter, Firebase, Cloud Firestore, and a
    TypeScript/Express API).
  \item It demonstrates a complete, real-money mobile payment integration
    (Monime, covering Orange Money and Afrimoney) embedded directly in the
    pickup-request workflow as a state-machine gate rather than a
    bolted-on afterthought, including handling for two undocumented
    provider-API behaviours discovered during integration
    (Section~\ref{sec:payment-module}).
  \item It documents, rather than hides, the specific points at which an
    ambitious original design (Google Maps Platform integration, a trained
    PyTorch classifier, a FastAPI backend) had to be adapted during
    implementation, together with the reasoning behind each substitution.
    This is offered as a realistic account of how a student software
    project's scope evolves under real technical and time constraints,
    consistent with the honesty requirement of NFR-ETH-9.
  \item It produces a substantial, reusable set of functional and
    non-functional requirements (Sections~\ref{sec:functional-requirements}
    and~\ref{sec:nfr}, Chapter~\ref{ch:analysis-design})
    and a module-by-module implementation and test-status audit
    (Table~\ref{tab:module-status}) that can guide any continuation of the
    project beyond this documentation.
\end{enumerate}

\section{Conclusion}

This study designed and implemented EcoTrace, a cross-platform electronic
waste recovery and circular economy management system, and evaluated it
against the ten specific objectives set out for it. The system's core
data-management, authentication, pickup-scheduling, collection, repair,
recycling, and payment functionality was implemented and, in the areas
covered by Section~\ref{sec:testing}, tested successfully, including one
genuine live-money mobile-money transaction. Two objectives -- automatic
AI-based classification using a trained model, and Google Maps Platform
integration -- were not delivered as originally specified; the first was
substituted with a pluggable heuristic classifier and the second with
device GPS and haversine-distance calculations, both of which are
functional but materially narrower than what was proposed. Formal usability
and field evaluation with real users, called for by the tenth objective,
was not performed within the scope of this study.

On balance, the study concludes that EcoTrace, as delivered, is a
functioning proof of the system envisioned in Chapter~\ref{ch:introduction}:
it demonstrates that a coordinated, cloud-based, cross-platform approach to
e-waste registration, collection, and payment is technically achievable in
the Sierra Leonean context, while also demonstrating -- through the
substitutions and outstanding items catalogued in
Table~\ref{tab:module-status} -- exactly which parts of that vision still
require further work before the system could be considered
production-ready.

\section{Recommendations}

Based on the outcome of this study, the following recommendations are made.

\subsection{Recommendations for Deployment and Operation}

\begin{enumerate}[leftmargin=1.4cm]
  \item The backend should be moved from Render's free tier to a paid tier
    or an equivalent dedicated hosting arrangement before any real-world
    pilot, to remove the cold-start delays and resource limits identified
    as a limitation in Section~\ref{sec:limitations-intro}.
  \item Production credentials for the notification providers (FCM, SMS,
    and email) referenced in Table~\ref{tab:module-status} should be
    provisioned and verified end-to-end before the system is relied upon
    for time-sensitive alerts such as pickup confirmations or payment
    status changes.
  \item Because the Monime webhook signature could not be independently
    verified during this study (Section~\ref{sec:payment-module}), the
    payment-confirmation logic that re-queries Monime's API for the
    authoritative transaction status, rather than trusting the webhook
    payload directly, should be retained even if signature verification
    later becomes possible, as defence in depth.
  \item Firestore security rules, and in particular any newly added
    collection, should be checked against the deny-by-default catch-all
    rule as a standard step before deployment; the missing
    \texttt{paymentTransactions} rule found and fixed during this study
    (T-05, Table~\ref{tab:test-cases}) shows how easily such a gap can
    otherwise reach production undetected.
\end{enumerate}

\subsection{Recommendations for Continued Development}

\begin{enumerate}[leftmargin=1.4cm]
  \item Replace the heuristic-fallback classifier with a trained image
    classification model once a sufficiently large and diverse labelled
    e-waste image dataset is available, using the pluggable classifier
    interface already implemented (Section~\ref{sec:classification-module})
    so that the rest of the system requires no change.
  \item Integrate a mapping and routing provider (whether Google Maps
    Platform as originally proposed, or an open alternative such as
    OpenStreetMap-based routing) to add interactive map display and visual
    route optimisation on top of the GPS/haversine foundation already in
    place.
  \item Complete the integration-test coverage flagged as outstanding for
    the majority of "Partial" modules in Table~\ref{tab:module-status};
    most of these modules are functionally complete and require test
    coverage rather than new development.
  \item Extend the Monime integration with support for insufficient-funds
    detection once Monime exposes a corresponding webhook event, to remove
    the current reliance on a fixed client-side timeout
    (Section~\ref{sec:payment-module}).
\end{enumerate}

\section{Suggestions for Further Work}
\label{sec:future-work}

The following are suggested as areas for further work beyond the scope of
this study:

\begin{enumerate}[leftmargin=1.4cm]
  \item \textbf{Formal usability and field evaluation.} A structured
    usability study, and ideally a small field pilot with real households,
    businesses, and collectors in a defined area of Sierra Leone, would
    directly address the tenth objective, which this study was not able to
    evaluate empirically.
  \item \textbf{Support for additional mobile-money providers and
    countries.} The current Monime integration covers Orange Money and
    Afrimoney within Sierra Leone; extending payment support to other
    providers or countries would require re-examining the
    \texttt{authorizedProviders} versus \texttt{authorizedPhoneNumber}
    behaviour documented in Section~\ref{sec:payment-module} for each new
    provider.
  \item \textbf{Offline support.} Given the Internet-dependency limitation
    identified in Section~\ref{sec:limitations-intro}, future work could
    add local caching and background synchronisation so that collectors
    and technicians in areas of unreliable connectivity can continue
    working and sync once connectivity is restored.
  \item \textbf{Environmental impact methodology validation.} The
    environmental-impact calculations implemented in the current system
    (Table~\ref{tab:module-status}) should be validated against an
    established environmental-accounting methodology or reviewed by a
    domain expert before being used for external reporting.
  \item \textbf{Independent security assessment.} An independent
    penetration test or security audit, beyond the functional and security
    testing reported in Section~\ref{sec:testing}, is recommended before
    any production deployment that will handle real payment traffic at
    scale.
\end{enumerate}

\section{Chapter Summary}

This chapter summarised the study, assessed the ten specific objectives set
out in Chapter~\ref{ch:introduction} against the implementation and testing
evidence reported in Chapter~\ref{ch:implementation}, and found that five
objectives were achieved, four were partially achieved, and one was not
achieved as originally specified, while a substantial payment-integration
module outside the original objectives was delivered and verified. It
concluded that EcoTrace, as implemented, is a working demonstration of the
system envisioned at the outset of the study, with specific, clearly
identified gaps remaining before it could be considered production-ready.
Recommendations were offered for deploying and operating the system as it
stands, for continuing its development, and for further work -- most
notably formal usability evaluation -- that falls outside the scope of this
study. This concludes the main body of the project documentation; the
references and appendices that follow provide supporting material.
